\documentclass[11pt,a4paper]{article}

% ── Packages ──────────────────────────────────────────────────────────────────
\usepackage[margin=1in]{geometry}
\usepackage{amsmath,amssymb}
\usepackage{enumitem}
\usepackage{listings}
\usepackage{xcolor}
\usepackage{hyperref}
\usepackage{titlesec}
\usepackage{booktabs}
\usepackage{mdframed}
\usepackage{parskip}
\usepackage{microtype}
\usepackage{fancyhdr}

% ── Colours ───────────────────────────────────────────────────────────────────
\definecolor{keycolor}{RGB}{0,100,180}
\definecolor{codebg}{RGB}{245,245,245}
\definecolor{codecomment}{RGB}{80,130,80}
\definecolor{codestring}{RGB}{180,50,50}
\definecolor{boxbg}{RGB}{240,248,255}
\definecolor{warnbg}{RGB}{255,248,230}
\definecolor{warnborder}{RGB}{200,140,0}

% ── Code listings ─────────────────────────────────────────────────────────────
\lstset{
  language=C,
  backgroundcolor=\color{codebg},
  basicstyle=\ttfamily\small,
  keywordstyle=\color{keycolor}\bfseries,
  commentstyle=\color{codecomment}\itshape,
  stringstyle=\color{codestring},
  numbers=left,
  numberstyle=\tiny\color{gray},
  numbersep=8pt,
  breaklines=true,
  frame=single,
  framerule=0.4pt,
  rulecolor=\color{gray!50},
  xleftmargin=2em,
  tabsize=4,
  showstringspaces=false,
}

% ── Callout boxes ─────────────────────────────────────────────────────────────
\newmdenv[
  backgroundcolor=boxbg,
  linecolor=keycolor,
  linewidth=1.5pt,
  leftline=true, rightline=false, topline=false, bottomline=false,
  innerleftmargin=10pt, innerrightmargin=8pt,
  skipabove=6pt, skipbelow=6pt,
]{keybox}

\newmdenv[
  backgroundcolor=warnbg,
  linecolor=warnborder,
  linewidth=1.5pt,
  leftline=true, rightline=false, topline=false, bottomline=false,
  innerleftmargin=10pt, innerrightmargin=8pt,
  skipabove=6pt, skipbelow=6pt,
]{warnbox}

% ── Section formatting ────────────────────────────────────────────────────────
\titleformat{\section}{\large\bfseries\color{keycolor}}{}{0em}{}[\titlerule]
\titleformat{\subsection}{\normalsize\bfseries}{}{0em}{}
\titleformat{\subsubsection}{\normalsize\itshape\bfseries}{}{0em}{}

% ── Header / footer ──────────────────────────────────────────────────────────
\pagestyle{fancy}
\fancyhf{}
\lhead{\small CS~3210 -- Operating Systems}
\rhead{\small Lectures 10 \& 11: Interrupts and Concurrency}
\cfoot{\small\thepage}
\renewcommand{\headrulewidth}{0.4pt}

% ── Macros ────────────────────────────────────────────────────────────────────
\newcommand{\term}[1]{\textbf{#1}}
\newcommand{\code}[1]{\texttt{#1}}
\newcommand{\xv}{xv6}

% ══════════════════════════════════════════════════════════════════════════════
\begin{document}

\begin{center}
  {\LARGE\bfseries CS~3210 -- Operating Systems}\\[4pt]
  {\large Lectures 10 \& 11: Interrupts, Exceptions, and Concurrency}\\[4pt]
  {\normalsize Alexey Tumanov \quad|\quad Georgia Tech SCS
    \quad|\quad Feb 17 \& 19, 2026}
\end{center}

\tableofcontents
\newpage

% ══════════════════════════════════════════════════════════════════════════════
\section{Virtual Memory: Swapping and the Clock Algorithm}
% ══════════════════════════════════════════════════════════════════════════════

\subsection{Memory Backing}

Physical RAM is \emph{not} the hard ceiling on usable memory.  Two mechanisms
extend the address space beyond physical RAM:

\begin{itemize}
  \item \term{Swapping} -- pages are evicted from RAM and saved to disk when
    memory is exhausted; transparently reloaded on access.
  \item \term{Memory-Mapped Files} -- file data is mapped directly into virtual
    addresses and written back to disk on modification, avoiding an explicit
    copy step.
\end{itemize}

\subsection{Page Eviction Policy}

When RAM is full, the OS must choose a victim page to evict.

\begin{center}
\begin{tabular}{lll}
  \toprule
  \textbf{Policy} & \textbf{Idea} & \textbf{Feasibility} \\
  \midrule
  Optimal (Bélády) & Evict the page used furthest in the future & Impossible (requires future knowledge) \\
  LRU              & Evict the least recently used page          & Theoretically ideal, practically costly \\
  Clock (LRU approx.) & Approximate LRU with a single access bit & Practical \\
  \bottomrule
\end{tabular}
\end{center}

\subsubsection{Why Pure LRU Is Expensive}

Exact LRU needs a timestamp on \emph{every} memory access (reads \emph{and}
writes).  The overheads are prohibitive:
\begin{itemize}
  \item Maintaining a sorted data structure costs
        $\mathcal{O}(n) \to \mathcal{O}(\log n) \to \mathcal{O}(1)$ depending
        on implementation, but \emph{all} require work on each access.
  \item The only practical hardware hook is a \term{page fault on every access},
        i.e.\ marking the PTE as not-present (\code{PTE\_P = 0}) and trapping
        to the OS -- this overhead is unacceptable even with infinite RAM.
\end{itemize}

\subsection{The Clock Algorithm (CAC)}

The clock algorithm approximates LRU using one \term{access bit}~\(a\) per
page, arranged in a circular list with a ``clock hand''.

\begin{keybox}
\textbf{Clock Algorithm Steps}
\begin{enumerate}[leftmargin=*]
  \item Hardware sets \(a = 1\) whenever a page is read or written.
  \item On eviction pressure, the clock hand scans pages in order:
    \begin{itemize}
      \item If \(a = 1\): clear to \(a = 0\) (``give a second chance'') and
            advance.
      \item If \(a = 0\): \textbf{evict} this page.
    \end{itemize}
  \item After eviction, advance the hand past the evicted page.
\end{enumerate}
\end{keybox}

\subsubsection{Trapping on Access}

The OS only needs to know when a page transitions from \emph{unaccessed} to
\emph{accessed}, not on every subsequent read.  The trick:

\begin{itemize}
  \item When the clock clears a page's access bit (\(a \leftarrow 0\)), it
        also clears \code{PTE\_P = 0}.
  \item The \emph{first} subsequent access faults, the fault handler sets
        \(a = 1\) and \code{PTE\_P = 1}, then resumes -- no further traps until
        the next clock sweep.
\end{itemize}

This means traps occur at most once per clock sweep per page, not on every
access.

\subsubsection{Worst-Case Analysis}

Worst case: all pages have \(a = 1\) when eviction is needed.  The clock hand
must do a full pass clearing all bits before it can evict -- $\mathcal{O}(n)$
in the number of pages.  In practice this is rare and the algorithm performs
well.

% ══════════════════════════════════════════════════════════════════════════════
\section{System Calls, Interrupts, and Exceptions}
% ══════════════════════════════════════════════════════════════════════════════

\subsection{Why Interrupts?}

Without interrupts, the only alternative is \term{polling}: the CPU repeatedly
checks whether an event has occurred.

\begin{itemize}
  \item \textbf{Infrequent events} (e.g.\ keyboard): polling wastes nearly all
        CPU cycles.
  \item \textbf{Frequent events}: polling may still miss events and couples
        event frequency to the polling rate.
\end{itemize}

Interrupts decouple the CPU from I/O timing entirely, improving both
\textbf{performance} and \textbf{correctness}.

\subsection{Interrupt Categories}

\begin{center}
\begin{tabular}{lll}
  \toprule
  \textbf{Category} & \textbf{Source} & \textbf{Examples} \\
  \midrule
  Exceptions       & CPU itself           & Page fault (vec.~14), divide-by-zero, memory errors \\
  Hardware (INTR)  & External devices     & Timer (vec.~32), keyboard, NIC, disk \\
  Software (INT)   & \code{INT n} instruction & System calls (\xv: \code{INT 0x40}) \\
  Debugging        & Breakpoints          & Software breakpoints \\
  NMI              & Critical hardware    & Power failure, memory corruption \\
  IPI              & Other processors     & Inter-processor interrupt \\
  \bottomrule
\end{tabular}
\end{center}

\subsection{Maskable vs.\ Non-Maskable Interrupts}

\begin{itemize}
  \item \term{INTR (Maskable)}: Ignored when the \code{IF} flag in \code{EFLAGS}
    is 0.  Controlled with:
    \begin{itemize}
      \item \code{sti} -- set interrupt flag (enable)
      \item \code{cli} -- clear interrupt flag (disable)
    \end{itemize}
  \item \term{NMI (Non-Maskable Interrupt)}: Always delivered; used for
    catastrophic events (power failure, memory corruption).  In \xv, an NMI
    causes an immediate kernel panic.  x86 assigns NMI to vector~2.
\end{itemize}

\subsection{Programmable Interrupt Controller (PIC 8259A)}

The x86 CPU has a limited number of physical interrupt pins.  The
\term{Programmable Interrupt Controller (8259A)} fans out many device lines
to a single CPU \code{INTR} pin, assigning each device a configurable vector
number.

\subsection{Interrupt Vectors and the IDT}

\begin{keybox}
Every interrupt/exception type is assigned an 8-bit \term{vector number}
(0--255).  When an interrupt fires, the vector indexes into the
\term{Interrupt Descriptor Table (IDT)} to find the handler address.
\end{keybox}

\subsubsection{Interrupt Descriptor Table (IDT)}

\begin{itemize}
  \item 256 entries $\times$ 8 bytes = 2\,KB table (analogous to the GDT).
  \item Each entry is a \term{gate descriptor}: segment selector + handler
    offset + privilege level (\code{DPL}).
  \item The CPU register \term{IDTR} holds the physical address and size of
    the current IDT; loaded with the \code{lidt} instruction.
  \item In \xv, the IDT is populated at boot in \code{trap.c} using the
    \code{SETGATE} macro.  Vectors jump to stubs in \code{vector.S}, which
    then jump to \code{alltraps} in \code{trapasm.S}.
\end{itemize}

\subsubsection{Notable Vector Assignments in xv6}

\begin{itemize}
  \item Vector 14 $\to$ page-fault handler (Lab~2)
  \item Vector 32 $\to$ timer/clock handler $\to$ scheduler (Lab~3)
  \item Vector 0x40 $\to$ system call handler
\end{itemize}

\subsection{Software Interrupts and System Calls}

The \code{INT n} instruction lets user-space intentionally invoke any handler:

\begin{lstlisting}[language={[x86masm]Assembler},numbers=none]
INT 0x40    ; invoke system call handler (xv6)
\end{lstlisting}

To return from an interrupt handler, use \code{iret}, which atomically
restores \code{CS}, \code{EIP}, and \code{EFLAGS} from the kernel stack.

\subsection{The \texttt{INT} Instruction: Step-by-Step}

When \code{INT 0x40} executes:
\begin{enumerate}
  \item Determine the vector number (0x40 here).
  \item Fetch the 8-byte IDT entry for vector 0x40 (base address from IDTR
        $+$ 0x40 $\times$ 8).
  \item Check \code{CPL $\leq$ DPL} in the descriptor (for software INT only).
  \item If the target has a lower privilege level than the current CPL, save
        \code{ESP}/\code{SS} to a CPU-internal register and reload them from
        the \term{Task State Segment (TSS)} -- this is the kernel stack switch.
  \item Push: user \code{SS}, user \code{ESP}, user \code{EFLAGS}, user
        \code{CS}, user \code{EIP}.
  \item Clear interrupt-enable and other \code{EFLAGS} bits.
  \item Load \code{CS}:\code{EIP} from the IDT descriptor -- now executing
        kernel handler code.
\end{enumerate}

\subsection{Interrupt Handling: High-Level View}

\begin{enumerate}
  \item \textbf{Identify} the interrupt (vector number).
  \item \textbf{Privilege check} -- is the CPL allowed to invoke this vector?
  \item \textbf{Stack switch} -- from user stack to dedicated kernel/handler
        stack.
  \item \textbf{Save registers} -- the \term{trap frame} captures all
        user-space registers.
  \item \textbf{Dispatch} -- jump to the appropriate kernel handler
        (\code{trap.c} in \xv).
  \item \textbf{Return} via \code{iret}, restoring saved state.
\end{enumerate}

In \xv:
\begin{itemize}
  \item \code{x86.h} defines the \code{trapframe} struct (registers pushed by
    hardware + software).
  \item \code{usys.S} defines system call stubs (issues \code{INT 0x40}).
  \item \code{trapasm.S} defines \code{alltraps} (saves registers, calls C
    handler) and \code{trapret} (restores and \code{iret}s).
  \item \code{trap.c} contains the C-level \code{trap()} dispatcher -- most
    of the code you will modify in labs.
\end{itemize}

% ══════════════════════════════════════════════════════════════════════════════
\section{Concurrency I: Interrupts as Pseudo-Concurrency}
% ══════════════════════════════════════════════════════════════════════════════

\begin{warnbox}
\textbf{Key Insight.} Interrupts create \emph{pseudo-concurrency}: an
interrupt handler can preempt kernel code at any instruction boundary, just
like a second thread.  Any shared data accessed by both kernel code and an
interrupt handler is a potential race condition.
\end{warnbox}

\subsection{The Lock + Interrupt Deadlock Problem}

Consider kernel code holding a spinlock when an interrupt fires, and the
interrupt handler also tries to acquire the same lock:

\begin{lstlisting}
// Kernel code              // Interrupt handler
Lock();                     Lock();       // DEADLOCK: already held!
  x = shared;                Shared = 7;
Unlock();                   Unlock();
\end{lstlisting}

The handler spins forever waiting for a lock held by the code it interrupted.

\subsection{Solution: Disable Interrupts Around Locks}

The safe pattern -- used in \xv:

\begin{lstlisting}
// Kernel code
Interrupt_disable();   // cli
  Acquire(lock);
    x = shared;
    Shared = 7;
  Release(lock);
Interrupt_enable();    // sti
\end{lstlisting}

\begin{keybox}
\textbf{\xv Rule:} Interrupts are disabled (via \code{pushcli}) whenever a
spinlock is acquired.  Mutual exclusion with interrupt handlers is achieved
through interrupt disabling, \emph{not} by having the handler skip locks.
\end{keybox}

% ══════════════════════════════════════════════════════════════════════════════
\section{Concurrency II: Mutual Exclusion and Locking}
% ══════════════════════════════════════════════════════════════════════════════

\subsection{The ``Too Much Milk'' Problem}

A classic thought experiment for mutual exclusion with shared memory:

\textbf{Scenario:} Alice and Bob share a fridge.  Each independently checks
if milk is present; if not, they go buy some.

\begin{itemize}
  \item \term{Shared variable}: presence of milk (boolean).
  \item \term{Safety}: at most one person buys milk.
  \item \term{Liveness}: if milk is needed, \emph{someone} buys it.
\end{itemize}

\textbf{The race:}
\begin{enumerate}
  \item Bob checks fridge $\to$ no milk.
  \item Alice checks fridge $\to$ no milk (before Bob returns).
  \item Both go buy milk $\to$ too much milk!
\end{enumerate}

The problem generalises to any \textbf{check-then-act} pattern on shared
state.

\subsection{Formalising the Problem}

\begin{itemize}
  \item ``Check fridge'' $\equiv$ \emph{read} a shared variable.
  \item ``Put milk in fridge'' $\equiv$ \emph{update} a shared variable.
  \item Bugs occur only for \emph{some} execution interleavings, making them
    hard to reproduce.
  \item Correct solutions must be safe for \emph{all} possible interleavings.
\end{itemize}

\subsection{Peterson's Algorithm (Theoretical Only)}

Peterson's algorithm achieves mutual exclusion between two threads using only
atomic stores and loads (flag variables + a \emph{turn} variable).

\begin{warnbox}
\textbf{Completely Impractical in Practice.}  Peterson's algorithm does not
scale beyond two threads, does not account for hardware/compiler reordering,
and has $\mathcal{O}(n)$ complexity variants for $n$ threads.  Mentioned here
for historical context only.  In real systems, hardware atomic primitives are
required.
\end{warnbox}

\subsection{C11 Atomic Primitives}

C11 (\code{<stdatomic.h>}) provides atomic operations with defined memory
ordering semantics:

\begin{center}
\begin{tabular}{ll}
  \toprule
  \textbf{Operation} & \textbf{Description} \\
  \midrule
  \code{atomic\_load()}            & Atomic read \\
  \code{atomic\_store()}           & Atomic write \\
  \code{atomic\_fetch\_add()}      & Atomic read-modify-write (add) \\
  \code{atomic\_exchange()}        & Atomic load + store in one op \\
  \code{atomic\_compare\_exchange()} & Conditional atomic load+store (CAS) \\
  \bottomrule
\end{tabular}
\end{center}

\subsection{Building a Spinlock with \texttt{atomic\_exchange}}

\begin{lstlisting}
// lock == 0: free;  lock == 1: held
int r = 1;
while (r != 0) {
    r = atomic_exchange(&lock, 1);  // atomically set lock=1, return old value
}
// critical section
atomic_store(&lock, 0);             // release
\end{lstlisting}

\begin{itemize}
  \item If \code{lock} was 0, \code{exchange} returns 0, loop exits --
        \textbf{lock acquired}.
  \item If \code{lock} was 1, \code{exchange} returns 1, loop continues --
        \textbf{spinning}.
  \item This is a \term{spinlock}: the thread busy-waits consuming CPU cycles.
\end{itemize}

% ══════════════════════════════════════════════════════════════════════════════
\section{Memory Models and Data Races}
% ══════════════════════════════════════════════════════════════════════════════

\subsection{The Sequential Consistency Assumption Is Wrong}

Programmers naturally assume each statement executes in order and is
immediately visible to all cores.  This is called \term{sequential
consistency (SC)}, but neither x86 nor C/C++ guarantee it by default.

\subsection{x86 Total Store Order (TSO)}

x86 uses the \term{Total Store Order (TSO)} memory model:

\begin{itemize}
  \item \textbf{Write--Write (W$\to$W)} ordering: respected -- stores appear
    globally in program order.
  \item \textbf{Read--Write (R$\to$W)} ordering: respected.
  \item \textbf{Write--Read (W$\to$R)} ordering: \emph{not} guaranteed --
    a store may not be visible before a subsequent load in the same thread.
\end{itemize}

\textbf{Concrete example:}

\begin{lstlisting}
// Initial: x = y = 0
// Thread T1              // Thread T2
movl $1, (x)             movl $1, (y)
movl (y), %eax           movl (x), %ebx
\end{lstlisting}

Under TSO, both \code{\%eax = 0} and \code{\%ebx = 0} is \emph{possible}
(both threads can observe stale values due to store-buffer reordering), even
though it appears impossible under sequential consistency.

\subsection{C11 DRF0 Semantics}

C/C++11 offers \term{DRF0} (Data-Race-Free 0) semantics:

\begin{keybox}
\textbf{DRF0 Guarantee:} If your program has \emph{no data races}, then it
behaves as if sequentially consistent.  If it \emph{does} have data races,
behaviour is \textbf{undefined}.
\end{keybox}

\textbf{Definition of a data race:} two or more concurrent accesses to the
same memory location, where at least one is a \emph{write}, with no
synchronisation between them.

\subsection{Practical Takeaway}

Write \term{data-race-free (DRF)} code:
\begin{itemize}
  \item Use locks, semaphores, condition variables, or barriers whenever
    shared data is accessed.
  \item Use C11 atomic primitives only in rare, well-understood cases
    (they have complicated semantics).
  \item When in doubt, \emph{overlock} (accept a performance hit) rather than
    \emph{underlock} (introduce subtle correctness bugs).
\end{itemize}

% ══════════════════════════════════════════════════════════════════════════════
\section{Deadlock and the Dining Philosophers}
% ══════════════════════════════════════════════════════════════════════════════

\subsection{Deadlock}

\term{Deadlock} occurs when a set of threads each hold a lock and wait for a
lock held by another -- a circular dependency with no progress.

A realistic \xv example: \code{ideintr} holds the IDE lock, then calls
\code{wakeup()}, which tries to acquire \code{ptable.lock}.  If another thread
holds \code{ptable.lock} and waits for the IDE lock, deadlock results.

\subsection{Dining Philosophers}

Five philosophers sit at a round table, each needing two forks (shared with
neighbours) to eat.  If all simultaneously pick up their \emph{left} fork and
wait for the right one, everyone starves -- \textbf{deadlock}.

\textbf{Solution -- Impose a Global Lock Ordering:}

\begin{keybox}
Assign each resource (fork) a global number.  \emph{Always} acquire locks in
\textbf{increasing numeric order}.  This breaks circular wait: if the last
philosopher acquires the lower-numbered fork first, the cycle is broken.
\end{keybox}

\xv uses at most two locks simultaneously (e.g.\ directory lock + file lock
in the filesystem) and enforces a consistent acquisition order.

% ══════════════════════════════════════════════════════════════════════════════
\section{The \texttt{volatile} Keyword}
% ══════════════════════════════════════════════════════════════════════════════

\begin{warnbox}
\textbf{\texttt{volatile} is widely misunderstood.}  It does \emph{not}:
\begin{itemize}
  \item Provide atomicity.
  \item Enforce memory ordering or fencing.
  \item Make a variable safe to share between threads.
\end{itemize}

\texttt{volatile} only tells the compiler that the variable may change
outside the current thread's execution -- preventing certain compiler
optimisations (e.g.\ caching in a register).

Its \emph{only} legitimate use in systems code is for \textbf{memory-mapped
device registers} -- not for thread synchronisation.  Use C11 atomic
primitives or proper locks instead.

\textit{Further reading:} \url{https://www.kernel.org/doc/html/latest/process/volatile-considered-harmful.html}
\end{warnbox}

% ══════════════════════════════════════════════════════════════════════════════
\section{Spinlocks in xv6}
% ══════════════════════════════════════════════════════════════════════════════

\subsection{The \texttt{spinlock} Struct}

\begin{lstlisting}
struct spinlock {
    atomic_uint locked;   // 0 = free, 1 = held
    // Debugging fields:
    char       *name;     // Lock name
    struct cpu *cpu;      // CPU holding the lock
    uint        pcs[10];  // Call stack (program counters) at acquisition
};
\end{lstlisting}

\subsection{\texttt{acquire()}}

\begin{lstlisting}
void acquire(struct spinlock *lk) {
    pushcli();   // Disable interrupts to prevent interrupt deadlock

    if (holding(lk))
        panic("acquire");  // Re-entrant acquisition would deadlock

    // Spin until the lock is free, then atomically grab it.
    // memory_order_acq_rel also acts as a memory fence:
    //   reads after acquire() are not reordered before it.
    while (atomic_exchange_explicit(&lk->locked, 1,
                                    memory_order_acq_rel) != 0)
        ;   // spin

    // Record debug info
    lk->cpu = cpu;
    getcallerpcs(&lk, lk->pcs);
}
\end{lstlisting}

\textbf{Key points:}
\begin{itemize}
  \item \code{pushcli()} increments a per-CPU interrupt-disable counter;
    \code{popcli()} decrements it and re-enables interrupts only when the
    counter reaches zero (supports nested acquire/release).
  \item \code{memory\_order\_acq\_rel} on the exchange is a \term{memory
    barrier}: ensures no memory operations from inside the critical section
    migrate before the acquire.
  \item Panics on re-entrant acquisition -- \xv spinlocks are non-recursive.
\end{itemize}

\subsection{\texttt{release()}}

\begin{lstlisting}
void release(struct spinlock *lk) {
    if (!holding(lk))
        panic("release");

    lk->pcs[0] = 0;
    lk->cpu = 0;

    // Release the lock with a release fence: ensures all writes inside the
    // critical section are visible before the lock is freed.
    atomic_store_explicit(&lk->locked, 0, memory_order_release);

    popcli();  // Re-enable interrupts (if outermost acquire)
}
\end{lstlisting}

\textbf{Key points:}
\begin{itemize}
  \item \code{memory\_order\_release} ensures critical-section writes are
    visible to the next thread that acquires the lock.
  \item \code{popcli()} after the store: interrupts must stay disabled until
    \emph{after} the lock is dropped to maintain the invariant.
\end{itemize}

% ══════════════════════════════════════════════════════════════════════════════
\section{Summary and Big Picture}
% ══════════════════════════════════════════════════════════════════════════════

\begin{center}
\begin{tabular}{ll}
  \toprule
  \textbf{Concept} & \textbf{One-Line Summary} \\
  \midrule
  Clock Algorithm        & Approximate LRU using a circular access-bit list \\
  Interrupts             & Asynchronous events handled via the IDT \\
  IDT / Vectors          & 256-entry table mapping event numbers to handlers \\
  INT instruction        & Software interrupt; used for syscalls in \xv \\
  Trap Frame             & Saved register state pushed on interrupt entry \\
  Interrupt Deadlock     & Handler + kernel code competing for the same lock \\
  pushcli / popcli       & Disable interrupts while holding a spinlock in \xv \\
  Too Much Milk          & Motivates atomic check-and-set primitives \\
  Peterson's Algorithm   & 2-thread software mutex -- impractical \\
  atomic\_exchange       & Hardware primitive for building spinlocks \\
  TSO Memory Model       & x86 allows W$\to$R reordering \\
  DRF0                   & No data race $\Rightarrow$ sequential consistency \\
  Dining Philosophers    & Deadlock; solution: global lock ordering \\
  volatile               & Compiler hint only; not a sync primitive \\
  xv6 spinlock           & atomic\_exchange + interrupt disable + debug info \\
  \bottomrule
\end{tabular}
\end{center}

\end{document}
